\section{Surrogate Fidelity and the Validation Gate}
\label{sec:hybrid_gate}

\subsection{Motivation}
An analytic optimum is only meaningful if the closed-form cost that produced it agrees with the high-fidelity DSMC solution at the same geometry. The cost of Eq.~\eqref{eq:ga_cost} is built from the Sutton--Graves correlation at the physical nose radius, the geometric drag law and the ballistic-coefficient ratio; it cannot represent particle noise, real-gas chemistry, or geometric discretization effects present in the DSMC solution. A correlation that is systematically optimistic would rank designs on an error rather than on physics. The validation gate closes this loop: it re-simulates a selected winner with the high-fidelity chain and measures the analytic-to-DSMC gap before the design is accepted.

\subsection{Gate Protocol}
The gate executes five legs as a post-optimization stage:
\begin{enumerate}
    \item \textbf{Preconditions and geometry QA.} The compiled binary, the Docker daemon, the validation virtual environment, the pipeline module, and the reference CSV are all checked, followed by a seconds-fast \texttt{\seqsplit{bin/main -{}-validate-only}} geometry check. Every failure is reported with a path and a remedy, and the gate exits with status 1.
    \item \textbf{DSMC leg.} The selected winner is re-simulated with \texttt{\seqsplit{bin/main -{}-validate -{}-headless}} at its geometry, writing to a dedicated results directory. The reference dataset used for the baseline comparison is never overwritten.
    \item \textbf{Refinement leg.} The validation pipeline is executed on the freshly written dataset, producing a new unified comparison dataset under the gate directory.
    \item \textbf{Verdict.} The gate passes if and only if $q_{hybrid} \leq Q_{target} = 165.716$\,J/cm$^2$, and it records the prediction error $(q_{hybrid} - q_{analytic})/q_{analytic}$.
    \item \textbf{Multi-fidelity correction.} A bounded correction factor $c = \operatorname{clamp}\bigl(q^{\mathrm{DSMC}}_{\mathrm{flux}}/\allowbreak q^{\mathrm{SG}}_{\mathrm{flux}},\,\allowbreak 0.5,\,\allowbreak 2.0\bigr)$ is formed from the run's own peak surface convection flux and the Sutton--Graves prediction for the same flown nose radius at the same flight state, and it is applied to the heat-load term of $J(\mathbf{x})$: the corrected objective is $J^{(2)}(\mathbf{x}) = J(\mathbf{x}) + (c-1)\,Q(\mathbf{x})/Q_{target}$, which the implementation realises as an effective heat-load weight of $c$. The same genetic algorithm is then re-run with that weight --- in closed form, with no solver call --- producing the corrected design of Sec.~\ref{sec:double_pass_ga}.
\end{enumerate}

\subsection{Final DSMC Verification}
The gate was exercised on the pass-1 optimum of this campaign: $R_N = 1.3359$\,m, $r_{tor} = 0.1796$\,m,
$\theta_c = 58.61^\circ$, physical nose radius $0.4898$\,m, six toroids, scalloped skin. The run used
2200 DSMC steps and the in-process PINN\_AUX leg, and its record is
\texttt{results/validation\_optimized}. It returned
\begin{equation}
    q_{hybrid} = 175.580\ \text{J/cm}^2 \;>\; Q_{target} = 165.716\ \text{J/cm}^2,
    \label{eq:gate_fail}
\end{equation}
so the pass-1 heat-load verdict is a \emph{fail}, by $51.88\,\% of the target. The analytic
screening model had predicted $364.87$\,J/cm\textsuperscript{2} for the same geometry, so the model
over-predicts the solver's heat load by $2.08$; the discrepancy is the quantity the correction
factor is built from.

The correction is formed inside that same run, from the run's own peak surface convection flux
($1\,688\,940$ W/m\textsuperscript{2}) and the Sutton--Graves prediction for the same flown nose radius
at the same flight state ($129\,293.3$ W/m\textsuperscript{2}). The raw ratio is $13.0629$, and
$c = \operatorname{clamp}(\cdot, 0.5, 2.0) = 2.0$ sits at the upper bound. The second pass
therefore runs the same genetic algorithm with an effective heat-load weight of $c = 2.0$, and the
resulting vehicle $R_N = 3.0000$\,m (physical $D = 6.0000$\,m, nose $1.1000$\,m, torus
$0.0500$\,m, half-cone $45.9656^\circ$) is re-simulated in
\texttt{results/validation\_optimized\_round2}. That run returns
\begin{equation}
    q_{hybrid} = 117.526\ \text{J/cm}^2 \;\leq\; Q_{target} = 165.716\ \text{J/cm}^2,
    \label{eq:gate_pass}
\end{equation}
so the pass-2 heat-load verdict is a \emph{pass}, with a margin of $\-29.08\,\%$. Its recorded
drag coefficient (unclamped) is $2.7813$ against the reference's $1.76551$, and its peak
deceleration is $8.357$ g against the reference's $22.720$ g.

Two properties of these runs matter for how their numbers are read. First, the drag coefficient in the
run's \texttt{cd} column is clamped to $[0,3]$ by the SPARTA writer: the pass-1 run saturates that
clamp on one of its twenty-two recorded steps and the pass-2 run on twelve, so the column is not flat in
either case, but any table that reads it must use the \texttt{cd\_unclamped} column instead. Second,
the correction factor is clamped, not measured to that value: the solver's peak surface flux is
$13.0629$ times the correlation's, and the clamp is what keeps the second pass from trusting a
single-point ratio without bound.

The reference dataset used for that comparison predates the extended CSV schema, so three of its
quantities are derived here rather than recorded. Its body-area drag coefficient follows from
\texttt{drag\_sum\_N}/(q\,A_{\mathrm{body}}) with $A_{\mathrm{body}} = \pi(3.0/2)^2 =
7.0686$\,m$^2$, and the derivation reproduces the recorded \texttt{cd} column to $10^{-6}$ at
both ends of the trajectory ($2.423278$ against $2.423279$ at step 100, and $1.765511$ against
$1.765511$ at step 2200), which also establishes that the reference never saturated the $[0,3]$ clamp.
Its DSMC surface heat load is $\max(\texttt{heatflux\_avg})\times 13.58$\,s $= 1933.13$\,J/cm$^2$.
Its flown geometry is the command-line default: diameter $3.0$\,m, nose $0.55$\,m, torus $0.135$\,m,
half-cone $60^\circ$, six toroids.

\subsection{Relation to the Current Optimum}
The gate was executed on the geometry of Table~\ref{tab:ga_before_after}, and its verdict on that geometry is a fail: the pass-1 candidate carries $175.580$\,J/cm\textsuperscript{{2}} against the $165.716$\,J/cm\textsuperscript{{2}} target. The design that the gate does admit is the pass-2 geometry of Table~\ref{tab:ga_two_pass}, a $6.0000$\,m diameter body, which passes the same gate at $117.526$\,J/cm\textsuperscript{{2}}. The pipeline therefore does what it is meant to do: the first analytic pass proposes a design the solver rejects, the solver's own flux ratio sets the correction, and the second analytic pass proposes a design the solver accepts. What the gate does not establish is a flight-qualified vehicle: every candidate is flown on the reference's matched trajectory, so the trajectory-geometry self-consistency of a $6.0000$\,m body on a $3.0$\,m body's flight path remains open, and it is stated as such rather than absorbed into the result.